\documentclass[a4paper, 11pt]{article}


% -- Coments
\usepackage{verbatim}

% -- Fuente --
% \usepackage{fontspec}
% \setmonofont{JetBrainsMono Nerd Font}  

% -- Color --
\usepackage{xcolor}
%\definecolor{azul}{RGB}{00,33,99}
\definecolor{azul}{RGB}{35,72,180}

% -- Page Margin --
\usepackage[margin=1in]{geometry}

% -- Espaciados --
\newcommand{\Pspace}{0.5cm}
\newcommand{\Aspace}{0.2cm}

% -- Columnas --
\usepackage{multicol}

% -- Imagenes --
\usepackage{graphicx}
\usepackage{float}

% -- Matemáticas --
\usepackage{amsmath, amssymb}

% -- Gráficas --
\usepackage{pgfplots}
\pgfplotsset{compat=1.18}

% -- Código --
% \usepackage{minted}

\usepackage{minted}
\setminted{
    fontsize=\small,
    breaklines=true,               % Ajusta el código al ancho de página
    frame=lines,                   % Añade líneas arriba y abajo
    framesep=2mm,
    baselinestretch=1.2,
    bgcolor=lightgray!10           % Color de fondo muy suave
}
\newcommand{\mips}[1]{\mintinline{mips}{#1}}
\newcommand{\asm}[1]{\mintinline{asm}{#1}}

\usepackage{listings}
\lstset{
    language=C++,                   % Lenguaje del código
    basicstyle=\ttfamily\small,     % Fuente del código
    keywordstyle=\color{blue},      % Color de palabras clave
    commentstyle=\color{gray},      % Color de comentarios
    stringstyle=\color{red},        % Color de cadenas
    numbers=left,                   % Números de línea a la izquierda
    numberstyle=\tiny\color{gray},
    breaklines=true,                % Permitir saltos de línea
    frame=single                    % Marco alrededor del código
}

\usepackage[colorlinks=true, linkcolor=blue, urlcolor=blue]{hyperref}

\begin{comment}
-- Formato de pregunta simple
    % - Problema n
    \vspace{\Pspace}
    \item Problema \par
        % Respuesta:
        \vspace{\Aspace} \par
        { \color{azul}  }


 -- Formato de pregunta multimple
    % - Problema 
    \vspace{\Pspace}
    \item  \par
        % Respuestas:
        \vspace{\Aspace} \par
        a) 
        \\ { \color{azul} }

        \vspace{\Aspace} \par
        b) 
        \\ { \color{azul} }

        \vspace{\Aspace} \par
        c) 
        \\ { \color{azul} }


 -- Formato de imagen
\begin{figure}[!ht]
    \centering
    \includegraphics[width=0.5\textwidth]{}
\end{figure}


 -- Formato de tabla
\begin{tabular}{c|c}
    \textbf{A} & \textbf{B} \\
    \hline
    Texto & Texto \\
    Texto & Texto
\end{tabular} \\
\end{comment}


\title
{
    Arquitectura de Computadoras 2026-1 \\
    Tarea 2
}

\begin{document}

    \maketitle

    \begin{center}
        \begin{tabular}{r|l}
            \textbf{Expediente} & \textbf{Nombre} \\ \hline
            223210350 & Amaya Soria Angel Alberto \\
            219208106 & Bórquez Guerrero Angel Fernando \\
            223215039 & Miranda Sánchez Javier Leonardo \\
            223203899 & Tostado Cortes Dante Alejandro
        \end{tabular}
    \end{center}

    \rule{\linewidth}{0.3mm}

    \vspace{0.3cm}
    \begin{enumerate}
    % - Problema 1
    \vspace{\Pspace}
    \item En el modelo de von Neumann, ¿qué problema se genera debido al uso de una memoria única para datos e instrucciones? \par
        % Respuesta:
        \vspace{\Aspace} \par
        { \color{azul} Cuello de botella. }



    % - Problema 2
    \vspace{\Pspace}
    \item Mencionar una técnica moderna que ayude a mitigar el problema del modelo de Von Neumann. \par
        % Respuesta:
        \vspace{\Aspace} \par
        { \color{azul} Arquitectura Harvard Modificada. Implementa dos memorias, una que almacena las instrucciones y otra que almacena los datos, lo que permite el acceso simultáneo. }



    % - Problema 3
    \vspace{\Pspace}
    \item Suponer una arquitectura de 32 bits y los siguientes tipos de datos:
    \begin{itemize}
        \item \mintinline{mips}{char} (1 byte)
        \item \mintinline{mips}{short} (2 bytes)
        \item \mintinline{mips}{int} (4 bytes)
    \end{itemize}
    Dada la siguiente estructura:
    \begin{minted}{mips}

        struct Dato
        {
            char a;
            int b;
            short c;
        }

    \end{minted}
        % Respuestas:
        \newpage
        a) Determinar el tamaño total de la estructura suponiendo que la memoria está alineada a 4 bytes.
        \\ { \color{azul} 

                \begin{tabular}{|c|c|c|c|c|c|c|c|c|c|c|c|c|}
                    \hline
                    \textbf{Dirección} & 0 & 1 & 2 & 3 & 4 & 5 & 6 & 7 & 8 & 9 & 10 & 11 \\
                    \hline
                    \textbf{Variable} & a &  &  &  & b & b & b & b & c & c &  & \\ 
                    \hline
                \end{tabular} \\\\ 
                Dado que b es de tipo int y ocupa 4 bytes entonces se recorre a la siguiente posición múltiplo de 4 dejando los bytes 1 al 3 como relleno, y dado que c es de tipo short con tamaño de 2 bytes esta puede ocupar la posición 8 al ser múltiplo de 2. Al final, se quedan de relleno los bytes 10 y 11 para que la estructura esté alineada a 4 bytes.
        }

        \vspace{\Aspace} \par
        b) Indicar dónde se insertan los bytes de relleno.
        \\ { \color{azul} Posiciones: 1, 2, 3, 10 y 11. }

        \vspace{\Aspace} \par
        c) ¿Cuál es la ventaja de tener la memoria alineada?
        \\ { \color{azul} Mejora el rendimiento al permitir que el procesador acceda a los datos en un solo ciclo de lectura. }



    % - Problema 4
    \vspace{\Pspace}
    \item Dado el siguiente número \mintinline{mips}{0xFEDCBA98} y suponiendo que se pide guardarlo en la dirección \mintinline{mips}{5000}. \par
        % Respuestas:
        \vspace{\Aspace} \par
        a) Representar cómo se guarda en la memoria si la CPU es big-endian.
        \\ { \color{azul} 
                Transformamos el número \mintinline{mips}{0xFEDCBA98} de base 16 a base 2. \\ 
                \begin{tabular}{|c|c|c|c|c|c|c|c|}
                    \hline
                    \textbf{F} & \textbf{E} & \textbf{D} & \textbf{C} & \textbf{B} & \textbf{A} & \textbf{9} & \textbf{8} \\
                    \hline
                    1111 & 1110 & 1101 & 1100 & 1011 & 1010 & 1001 & 1000 \\
                    \hline
                    \multicolumn{2}{|c|}{Byte más alto} & \multicolumn{4}{c}{} & \multicolumn{2}{|c|}{Byte más bajo} \\
                    \hline
                \end{tabular} \\

                Representación Big-endian\\
                \begin{tabular}{|c|cc|}
                    \hline
                    \textbf{Dirección} & \multicolumn{2}{|c|}{\textbf{Byte}} \\
                    \hline
                    5000 & 1111 & 1110 \\
                    5001 & 1101 & 1100 \\
                    5002 & 1011 & 1010 \\
                    5003 & 1001 & 1000 \\
                    \hline
                \end{tabular} \\
        }

        \vspace{\Aspace} \par
        b) Representar cómo se guarda en la memoria si la CPU es little-endian.
        \\ { \color{azul} 
                Representación Little-endian\\
                \begin{tabular}{|c|cc|}
                    \hline
                    \textbf{Dirección} & \multicolumn{2}{|c|}{\textbf{Byte}} \\
                    \hline
                    5000 & 1001 & 1000 \\
                    5001 & 1011 & 1010 \\
                    5002 & 1101 & 1100 \\
                    5003 & 1111 & 1110 \\
                    \hline
                \end{tabular} \\
        }

        \vspace{\Aspace} \par
        c) Si en ambos casos, se lee un \mintinline{mips}{short} en la dirección \mintinline{mips}{5000}, ¿qué valor se obtiene?
        \\ { \color{azul} 
            Big-endian: 1111 1110 1101 1100 \\
            Little-endian: 1011 1010 1001 1000
        }



    % - Problema 5
    \newpage
    \item Suponer una ISA con instrucciones de tamaño fijo de 4 bytes y 64 registros: \par
        % Respuestas:
        \vspace{\Aspace} \par
        a) ¿Cuántos bits se requieren para especificar un registro?
        \\ { \color{azul} 
            Dado que cada registro necesita un número único, entonces se requiere una cantidad de bits de $\log_{2}(64)$, lo que equivale a una longitud de 6 bits.
        }

        \vspace{\Aspace} \par
        b) Si una instrucción tipo R necesita 3 registros, ¿cuántos bits quedan para el opcode y campos adicionales?
        \\ { \color{azul} 
            3 registros equivalen a 18 bits, y dado que el tamaño de nuestras instrucciones es de 4 bytes que son 32 bits, entonces nos quedan $32 - 18 = 14$ bits para el opcode y campos adicionales.
        }

        \vspace{\Aspace} \par
        c) ¿Cuál es la ventaja de usar longitud fija?
        \\ { \color{azul} 
            Como el procesador sabe que cada instrucción mide lo mismo, puede empezar a buscar la siguiente instrucción antes de terminar la actual.
        }



    % - Problema 6
    \vspace{\Pspace}
    \item Dado el siguiente código en C:
    \begin{minted}{c}
        a = b + c;
        d = a + c;
    \end{minted}
    Escribir la secuencia de instrucciones en:
        % Respuestas:
        \vspace{\Aspace} \par
        a) Arquitectura de pila.
        { \color{azul} 
            \begin{minted}{mips}
        PUSH b
        PUSH c
        ADD
        POP a
        PUSH a
        PUSH c
        ADD
        POP d
            \end{minted}
        }

        \vspace{\Aspace} \par
        b) Arquitectura de acumulador.
        { \color{azul} 
            \begin{minted}{mips}
        LOAD b
        ADD c
        STORE a
        ADD c
        STORE d
            \end{minted}
        }

        \vspace{\Aspace} \par
        c) Arquitectura GPR memoria a memoria (tipo x86).
        { \color{azul} 
            \begin{minted}{mips}
        ADD a, b, c
        ADD d, a, c
            \end{minted}
        }



    % - Problema 7
    \vspace{\Pspace}
    \item ¿Qué restricción impone la arquitectura load/store? \par
        % Respuesta:
        \vspace{\Aspace} \par
        { \color{azul} Solo las instrucciones load y store pueden interactuar con la memoria. }



    % - Problema 8
    \vspace{\Pspace}
    \item Suponer una CPU RISC con instrucciones de tamaño fijo de 4 bytes y una CPU CISC con instrucciones de tamaño variable y un promedio de 18 bits por instrucción. Suponer que un programa se compila en la CPU RISC y produce 120 instrucciones y que luego se compila en la CPU CISC y produce 95 instrucciones. \par
        % Respuestas:
        \vspace{\Aspace} \par
        a) ¿Cuál es el tamaño total del programa en ambas arquitecturas?
        \\ { \color{azul} 
            RISC: 4 bytes = 32 bits $*$ 120 = 3840 bits \\
            Para CISC dado que tenemos un tamaño variable promedio entonces su tamaño total es aproximado. \\
            CISC: 18 bits $*$ 95 $\approx$ 1710 bits
        }

        \vspace{\Aspace} \par
        b) ¿Cuál tiene mayor densidad?
        \\ { \color{azul} 
            Como el programa en CISC logra realizar la misma tarea ocupando menos espacio total en memoria entonces este tiene una mayor densidad.
        }

        \vspace{\Aspace} \par
        c) ¿Cuál es la ventaja de usar instrucciones de tamaño variable?
        \\ { \color{azul} 
            La ventaja es en el uso en memoria. Al permitir el tamaño variable las instrucciones pueden ir desde 1 hasta 15 bytes aprovechando mejor el espacio al utilizar solo la cantidad de bytes necesarios.
        }



    % - Problema 9
    \vspace{\Pspace}
    \item Suponer una ISA load/store con instrucciones de tamaño fijo de 32 bits y que las instrucciones de carga y store utilizan 8 bits para el opcode, 5 bits para el registro base y el resto para el offset. Determinar cual es el rango de direcciones que se puede accesar, suponiendo que el offset es un entero con signo. \par
        % Respuesta:
        \vspace{\Aspace} \par
        { \color{azul}  
            \textbf{Calculamos los bits disponibles para el offset} \\
            $32 \text{(total)} - 8 \text{(opcode)} - 5 \text{(registro)} = 19$ \\
    
            \textbf{Calculamos el rango de direcciones} \\
            Dado que el offset es un entero con signo sus direcciones van desde $-2^{n-1}$ hasta $2^{n-1} - 1$: \\
            $-2^{19-1} = -2^{18} = -262144$ \\
            $2^{19-1} - 1 = 2^{18} - 1 = 262144 - 1 = 262143$ \\
            R/: El rango de direcciones del offset es -262144 hasta 262143.
        }



    % - Problema 10
    \vspace{\Pspace}
    \item Dado el número \mintinline{mips}{0x8D09FFFC} indicar qué instrucción MIPS es. \par
        % Respuesta:
        \vspace{\Aspace}
        { \color{azul}  
            \textbf{Transformamos el número \mintinline{mips}{0x8D09FFFC} de base 16 a base 2} \\ 
            \begin{tabular}{|c|c|c|c|c|c|c|c|}
                \hline
                \textbf{8} & \textbf{D} & \textbf{0} & \textbf{9} & \textbf{F} & \textbf{F} & \textbf{F} & \textbf{C} \\
                \hline
                1000 & 1101 & 0000 & 1001 & 1111 & 1111 & 1111 & 1100 \\
                \hline
            \end{tabular} \\

            \textbf{Descomponemos en el formato MIPS tipo I} \\
            \begin{tabular}{|c|c|c|c|}
                \hline
                \textbf{Opcode (6 bits)} & \textbf{rs (5 bits)} & \textbf{rt (5 bits)} & \textbf{Immediate/Offset (16 bits)} \\
                \hline 
                10 0011 & 0 1000 & 0 1001 & 1111 1111 1111 1100 \\
                \hline
                lw & \$t0 & \$t1 & -4 \\
                \hline
            \end{tabular} \\\\
            Instrucción: \mintinline{mips}{lw $t1, -4($t0)}
        }



    % - Problema 11
    \vspace{\Pspace}
    \item Dado el número \mintinline{mips}{0x8D2A0008} indicar qué instrucción MIPS es y cuál es la dirección efectiva. \par
        % Respuesta:
        \vspace{\Aspace} \par
        { \color{azul}  
            \textbf{Transformamos el número \mintinline{mips}{0x8D2A0008} de base 16 a base 2} \\ 
            \begin{tabular}{|c|c|c|c|c|c|c|c|}
                \hline
                \textbf{8} & \textbf{D} & \textbf{2} & \textbf{A} & \textbf{0} & \textbf{0} & \textbf{0} & \textbf{8} \\
                \hline
                1000 & 1101 & 0010 & 1010 & 0000 & 0000 & 0000 & 1000 \\
                \hline
            \end{tabular} \\

            \textbf{Descomponemos en el formato MIPS tipo I} \\
            \begin{tabular}{|c|c|c|c|}
                \hline
                \textbf{Opcode (6 bits)} & \textbf{rs (5 bits)} & \textbf{rt (5 bits)} & \textbf{Immediate/Offset (16 bits)} \\
                \hline 
                10 0011 & 0 1001 & 0 1010 & 0000 0000 0000 1000 \\
                \hline
                lw & \$t1 & \$t2 & 8 \\
                \hline
            \end{tabular} \\\\
            Instrucción: \mintinline{mips}{lw $t2, 8($t1)} \\
            Dirección efectiva: \mintinline{mips}{$t1 + 8}
        }



    % - Problema 12
    \vspace{\Pspace}
    \item Dado el número \mintinline{mips}{0x0800002A} indicar qué instrucción MIPS es y cuál es al dirección efectiva. \par
        % Respuesta:
        \vspace{\Aspace} \par
        { \color{azul}  
            \textbf{Transformamos el número \mintinline{mips}{0x0800002A} de base 16 a base 2} \\ 
            \begin{tabular}{|c|c|c|c|c|c|c|c|}
                \hline
                \textbf{0} & \textbf{8} & \textbf{0} & \textbf{0} & \textbf{0} & \textbf{0} & \textbf{2} & \textbf{A} \\
                \hline
                0000 & 1000 & 0000 & 0000 & 0000 & 0000 & 0010 & 1010 \\
                \hline
            \end{tabular} \\

            \textbf{Descomponemos en el formato MIPS tipo J} \\
            \begin{tabular}{|c|c|}
                \hline
                \textbf{Opcode (6 bits)} & \textbf{Address (26 bits)} \\
                \hline 
                00 0010 & 00 0000 0000 0000 0000 0010 1010 \\
                \hline
                j & 42 \\
                \hline
            \end{tabular} \\\\
            Instrucción: \mintinline{mips}{j 42} \\
            Dirección efectiva: \mintinline{mips}{0xA8}
        }



    % - Problema 13
    \vspace{\Pspace}
    \item Dado el siguiente código MIPS:
    \begin{minted}{mips}
        .data
            A: .word 1, 2, 3, 4, 5
            B: .word 21, 22, 23, 24, 25
        .text
            la $t0, B
            lw $t1, -8($t0)
    \end{minted}
    Escribir los valores finales de \mintinline{mips}{$t0} y \mintinline{mips}{$t1}. \par
        % Respuesta:
        \vspace{\Aspace} \par
        { \color{azul}  
            Creemos una tabla con direcciones relativas \\\\
            \begin{tabular}{|c|c|c|}
                \hline
                \textbf{Dirección relativa} & \textbf{Etiqueta} & \textbf{Valor} \\
                \hline
                0x10010000 & A & 1 \\
                0x10010004 & A & 2 \\
                0x10010008 & A & 3 \\
                0x1001000C & A & 4 \\
                0x10010010 & A & 5 \\
                \hline
            \end{tabular}
            \begin{tabular}{|c|c|c|}
                \hline
                \textbf{Dirección relativa} & \textbf{Etiqueta} & \textbf{Valor} \\
                \hline
                0x10010014 & B & 21 \\
                0x10010018 & B & 22 \\
                0x1001001C & B & 23 \\
                0x10010020 & B & 24 \\
                0x10010024 & B & 25 \\
                \hline
            \end{tabular} \\\\
            $\rightarrow$ "\mintinline{mips}{la $t0, B}" guarda el la dirección de \mintinline{mips}{B} en \mintinline{mips}{$t0}, siendo esta \mintinline{mips}{0x10010014}. \\
            $\rightarrow$ "\mintinline{mips}{lw $t1, -8($t0)}" guarda en \mintinline{mips}{$t1} el valor de la dirección \mintinline{mips}{B} -8 bytes, siendo este la dirección \mintinline{mips}{0x1001000C} con un valor de 4. \\
            R/: \mintinline{mips}{$t0 = 0x10010014} y \mintinline{mips}{$t1 = 4}.
        }



    % - Problema 14
    \vspace{\Pspace}
    \item Dado el siguiente código MIPS:
    \begin{minted}{mips}
        .data
            A: .word 11, 12, 13, 14, 15
        .text
            la $t0, A
            lw $t1, 0($t0)
            lw $t2, 12($t0)
            mul $t3, $t1, $t2
    \end{minted}
    Escribir los valores finales de \mintinline{mips}{$t0}, \mintinline{mips}{$t1}, \mintinline{mips}{$t2} y \mintinline{mips}{$t3} \par
        % Respuesta:
        \vspace{\Aspace} \par
        { \color{azul}  
            Creemos una tabla con direcciones relativas \\\\
            \begin{tabular}{|c|c|c|}
                \hline
                \textbf{Dirección relativa} & \textbf{Etiqueta} & \textbf{Valor} \\
                \hline
                0x10010000 & A & 11 \\
                0x10010004 & A & 12 \\
                0x10010008 & A & 13 \\
                0x1001000C & A & 14 \\
                0x10010010 & A & 15 \\
                \hline
            \end{tabular} \\\\
            $\rightarrow$ "\mips{la $t0, A}" carga la dirección \mips{0x10010000} en \mips{$t0}. \\
            $\rightarrow$ "\mips{lw $t1, 0($t0)}" guarda el valor 11 en \mips{$t1}. \\
            $\rightarrow$ "\mips{lw $t2, 12($t0)}" guarda el valor 14 en \mips{$t2}. \\
            $\rightarrow$ "\mips{mul $t3, $t1, $t2}" multiplica los valores de \mips{$t1} y \mips{$t2} y guarda el resultado en \mips{$t3}, siendo este resultado 154. \\
            R/: \mips{$t0 = 0x10010000}, \mips{$t1 = 11}, \mips{$t2 = 14} y \mips{$t3 = 154}.
        }



    % - Problema 15
    \vspace{\Pspace}
    \item Suponer que \mintinline{mips}{$s0} vale \mintinline{mips}{0x80000000} y que \mintinline{mips}{$s1} vale \mintinline{mips}{0xD0000000}, para las siguientes instrucciones MIPS indicar el valor de los registros destino y si ese valor es el deseado o si ocurrió overflow. \par
        % Respuestas:
        \vspace{\Aspace} \par
        a) \mintinline{mips}{add $t0, $s0, $s1}
        \\ { \color{azul} 
            \mips{$s0 = 0x80000000} \\
            \mips{$s1 = 0xD0000000} \par
            \begin{tabular}{r|l}
                \mips{0x80000000} & \\
                \mips{0xD0000000} & + \\
                \hline
                \mips{0x150000000} & \\
            \end{tabular} \par
            Como el registro es de únicamente 32 bits el "1" de la izquierda se lleva de carry, entonces desaparece y solamente queda \mips{0x50000000}. \par
            \mips{$t0 = 0x50000000} \par
            Siguiendo las reglas donde \textbf{Negativo + Negativo = Postivo} significa un overflow y nuestro resultado final fue un número positivo, concluimos que hubo un overflow.
        }

        \newpage
        b) \mintinline{mips}{sub $t1, $s0, $s1}
        \\ { \color{azul}
            \mips{$s0 = 0x80000000} \\
            \mips{$s1 = 0xD0000000} \\
            Dado que \mips{0xD0000000} es complemento a 2 podemos reemplazar la resta por una suma invirtiendo los bits de \mips{$s1}. \par
            Complemento de \mips{0xD0000000}: \mips{0x30000000}

            \begin{tabular}{r|l}
                \mips{0x80000000} & \\
                \mips{0x30000000} & + \\
                \hline
                \mips{0xB0000000} & \\
            \end{tabular} \par
            \mips{$t1 = 0xB0000000} \par
            Se ha obtenido el valor deseado.
        }

        \vspace{\Aspace} \par
        c) \mintinline{mips}{add $t0, $s0, $s1} y luego \mintinline{mips}{add $t0, $t0, $s0}
        \\ { \color{azul} 
            \mips{$s0 = 0x80000000} \\
            \mips{$s1 = 0xD0000000} \\
            \begin{tabular}{r|l}
                \mips{0x80000000} & \\
                \mips{0xD0000000} & + \\
                \hline
                \mips{0x150000000} & \\
            \end{tabular} \par
            \mips{$t0 = 0x50000000} \par
            \begin{tabular}{r|l}
                \mips{0x50000000} & \\
                \mips{0x80000000} & + \\
                \hline
                \mips{0xD0000000} & \\
            \end{tabular} \par
            \mips{$t0 = 0xD0000000} \par
            Aunque la segunda instrucción \mips{add} no provoca un overflow por sí misma (al sumar un positivo con un negativo), la primera operación sí generó un desbordamiento. Por lo tanto, el valor final en \mips{$t0} es incorrecto desde el punto de vista de la aritmética de números con signo.
        }



    % - Problema 16
    \vspace{\Pspace}
    \item Suponer que \mintinline{mips}{$t0} vale \mintinline{mips}{0xAAAAAAAA} y que \mintinline{mips}{$t1} vale \mintinline{mips}{0x12345678}. Para las siguientes instrucciones MIPS indicar el valor de los registros destino\par
        % Respuestas:
        \vspace{\Aspace} \par
        a) \mintinline{mips}{sll $t2, $t0, 4} y luego \mintinline{mips}{or $t2, $t2, $t1}.
        \\ { \color{azul} 
            \textbf{Shift Left Logical (\mips{sll})} \par
            \mips{$t0 = 0xAAAAAAAA} \par
            \mips{$t1 = 0x12345678} \par
            $\rightarrow$ Transformamos \mips{$t0} en binario. \par
            \mips{$t0 = 1010 1010 1010 1010 1010 1010 1010 1010} \par
            $\rightarrow$ Aplicamos \mips{sll} con 4 bits. \par
            \begin{tabular}{r|l}
                1010 1010 1010 1010 1010 1010 1010 1010 & $\leftarrow$ 4 bits \\
                \hline
                1010 1010 1010 1010 1010 1010 1010 0000 &
            \end{tabular} \par
            \mips{$t2 = 1010 1010 1010 1010 1010 1010 1010 0000} \par

            \newpage
            \textbf{Or (\mips{or})} \par
            \begin{tabular}{|c|c|c|c|}
                \hline
                \mips{$t2} & \mips{$t1} & \mips{or} & \textbf{Hex} \\
                \hline
                1010 & 0001 & 1011 & B \\
                1010 & 0010 & 1010 & A \\
                1010 & 0011 & 1011 & B \\
                1010 & 0100 & 1110 & E \\
                1010 & 0101 & 1111 & F \\
                1010 & 0110 & 1110 & E \\
                1010 & 0111 & 1111 & F \\
                0000 & 1000 & 1000 & 8 \\
                \hline
            \end{tabular} \par
            \mips{$t2 = BABEFEF8}
        }

        \vspace{\Aspace} \par
        b) \mintinline{mips}{sll $t3, $t0, 4} y luego \mintinline{mips}{andi $t3, $t3, -1}.
        \\ { \color{azul} 
            \textbf{Shift Left Logical (\mips{sll})} \par
            \mips{$t0 = 0xAAAAAAAA} \par
            \mips{$t1 = 0x12345678} \par
            $\rightarrow$ Transformamos \mips{$t0} en binario. \par
            \mips{$t0 = 1010 1010 1010 1010 1010 1010 1010 1010} \par
            $\rightarrow$ Aplicamos \mips{sll} con 4 bits. \par
            \begin{tabular}{r|l}
                1010 1010 1010 1010 1010 1010 1010 1010 & $\leftarrow$ 4 bits \\
                \hline
                1010 1010 1010 1010 1010 1010 1010 0000 &
            \end{tabular} \par
            \mips{$t3 = 1010 1010 1010 1010 1010 1010 1010 0000} \par

            \vspace{\Aspace}
            \textbf{And Immediate (\mips{andi})} \par
            \begin{tabular}{|c|c|c|c|}
                \hline
                \mips{$t3} & \mips{-1} & \mips{andi} & \textbf{Hex} \\
                \hline
                1010 & 1111 & 1010 & A \\
                1010 & 1111 & 1010 & A \\
                1010 & 1111 & 1010 & A \\
                1010 & 1111 & 1010 & A \\
                1010 & 1111 & 1010 & A \\
                1010 & 1111 & 1010 & A \\
                1010 & 1111 & 1010 & A \\
                0000 & 1111 & 0000 & 0 \\
                \hline
            \end{tabular} \par
            \mips{$t3 = AAAAAAA0}
        }

        \vspace{\Aspace} \par
        c) \mintinline{mips}{srl $t4, $t0, 3} y luego \mintinline{mips}{andi $t4, $t4, 0xFFEF}.
        \\ { \color{azul} 
            \textbf{Shift Right Logical (\mips{srl})} \par
            \mips{$t0 = 0xAAAAAAAA} \par
            \mips{$t1 = 0x12345678} \par
            $\rightarrow$ Transformamos \mips{$t0} en binario. \par
            \mips{$t0 = 1010 1010 1010 1010 1010 1010 1010 1010} \par
            $\rightarrow$ Aplicamos \mips{srl} con 3 bits. \par
            \begin{tabular}{r|l}
                1010 1010 1010 1010 1010 1010 1010 1010 & $\rightarrow$ 3 bits \\
                \hline
                0001 0101 0101 0101 0101 0101 0101 0101 &
            \end{tabular} \par
            \mips{$t4 = 0001 0101 0101 0101 0101 0101 0101 0101} \par

            \newpage
            \vspace{\Aspace}
            \textbf{And Immediate (\mips{andi})} \par
            $\rightarrow$ El valor inmediato \mips{0xFFEF} es de 16 bits. En operaciones lógicas inmediatas, MIPS realiza una extensión de ceros, por lo que el valor real de comparación es \mips{0x0000FFEF}. \par
            \begin{tabular}{|c|c|c|c|}
                \hline
                \mips{$t4} & \mips{0xFFEF} & \mips{andi} & \textbf{Hex} \\
                \hline
                0001 & 0000 & 0000 & 0 \\
                0101 & 0000 & 0000 & 0 \\
                0101 & 0000 & 0000 & 0 \\
                0101 & 0000 & 0000 & 0 \\
                0101 & 1111 & 0101 & 5 \\
                0101 & 1111 & 0101 & 5 \\
                0101 & 1110 & 0100 & 4 \\
                0101 & 1111 & 0101 & 5 \\
                \hline
            \end{tabular} \par
            \mips{$t4 = 00005545}
        }



    % - Problema 17
    \vspace{\Pspace}
    \item ARM tiene esta instrucción:
    \begin{minted}{asm}
        ADDNE r1, r2, r3
    \end{minted}
    Investigar qué hace y cuál es su ventaja sobre un branch. \par
        % Respuesta:
        \vspace{\Aspace} \par
        { \color{azul}  
            En ARM la mayoría de instrucciones tienen un sufijo condicional, en este caso \mips{ne}. \par
            
            \mips{add}: Operación \textbf{suma}. \\
            \mips{ne}: \textbf{n}ot \textbf{e}qual, que significa "no igual". \par

            La instrucción completa es \mips{r1 = r2 + r3} \textbf{solo si} el flag de estado indica que la comparación anterior resultó en \textbf{no igual}. \par

            \textbf{Por ejemplo} \\
            Supongamos que queremos sumarle puntos a un equipo de fútbol americano, estos puntos pueden ser 0, 3 o 7, pero únicamente queremos sumar puntos cuando no sean 0. \par
            $\rightarrow$ Código MIPS \\
            \mips{$t0}: Puntos totales \\
            \mips{$t1}: Puntos en la última jugada
            \begin{minted}{mips}
        ...
        beq $t1, $zero, omitir_suma
        add $t0, $t0, $t1
        j continuar
        omitir_suma:
            ...
        continuar:
            ...
            \end{minted} 
        
            \newpage
            $\rightarrow$ Código ARM \par
            \mintinline{asm}{r0}: Puntos totales \\
            \mintinline{asm}{r1}: Puntos en la última jugada
            \begin{minted}{asm}
        ...
        CMP r1, #0
        ADDNE r0, r0, r1
            \end{minted}
            \textbf{Ventaja} \\
            Los procesadores modernos usan \textbf{pipelining}, ejecutando varias etapas de distintas instrucciones simultáneamente. Cuando hay un salto de control (branch), el procesador debe predecir el camino a seguir; si la predicción falla, se debe vaciar el pipeline, perdiendo varios ciclos de reloj. Al utilizar \mintinline{asm}{ADDNE}, ARM elimina la necesidad de saltos condicionales para tareas simples, manteniendo el flujo de ejecución lineal. Esto resulta en un pipeline más eficiente, menor consumo de energía y una mayor densidad de código.
        }



    % - Problema 18
    \vspace{\Pspace}
    \item Dado el siguiente código en x86 (sintaxis intel): \par
    \begin{minted}{asm}
        mov eax, [A]
        add eax, [B]
        mov [C], eax
    \end{minted}
    Indicar a qué operación de alto nivel corresponde y cuál es la diferencia fundamental con un lenguaje RISC tipo MIPS.
        % Respuesta:
        \vspace{\Aspace} \par
        { \color{azul}  
            \textbf{Operación de alto nivel} \\
            $\rightarrow$ \asm{mov eax, [A]} \\
            Carga el contendio de \asm{A} en el registro \asm{eax}. \\
            
            $\rightarrow$ \asm{add eax, [B]} \\
            Suma el valor contendio en \asm{eax} con el valor contendio en \asm{B} y lo guarda en el registro \asm{eax}. \\
                
            $\rightarrow$ \asm{mov [C], eax} \\
            Guarda el resultado en la dirección \asm{C}. \\

            $\rightarrow$ Operación \\
            \mips{C = A + B} \par

            \textbf{Diferencia fundamental} \\
            La diferencia reside en las filosofías de diseño CISC (x86) frente a RISC (MIPS). \\
            x86 permite operandos de memoria directamente en instrucciones aritméticas (como en \asm{add eax, [B]}). En cambio, MIPS es una arquitectura \textbf{Load/Store}, lo que obliga a separar el acceso a memoria (\mips{lw}/\mips{sw}) de la operación lógica o aritmética. \\
            Mientras que en RISC cada instrucción realiza una sola tarea atómica para facilitar el pipelining, en CISC una sola instrucción puede realizar el ciclo completo de "Lectura-Modificación-Escritura", reduciendo el número total de líneas de código a costa de una mayor complejidad en el hardware del procesador.
        }



    % - Problema 19
    \vspace{\Pspace}
    \item En x86, ¿cuál es la diferencia entre la bandera de carry y la bandera de overflow? ¿Cuándo se prende o se apaga cada una de ellas? \par
        % Respuesta:
        \vspace{\Aspace} \par
        { \color{azul}  
            La diferencia fundamental radica en cómo se interpretan los datos: la bandera de \textbf{Carry} se utiliza para aritmética sin signo, mientras que la de \textbf{Overflow} se utiliza para aritmética con signo.

            \begin{itemize}
                \item \textbf{Carry Flag (CF):} Se activa cuando el resultado de una operación requiere un bit adicional de transporte fuera del bit más significativo (MSB). En una resta, actúa como un ``préstamo". Indica que el valor está fuera del rango $[0, 2^n - 1]$.
                \item \textbf{Overflow Flag (OF):} Se activa cuando el resultado de una operación con signo es erróneo porque excede el rango representable en complemento a 2..
            \end{itemize}

            \textbf{Criterios de activación:}
            \begin{itemize}
                \item \textbf{CF:} Se prende si hay un acarreo saliente del bit más significativo durante una suma, o un préstamo en una resta.
                \item \textbf{OF:} Se prende si el bit de acarreo que entra al MSB es diferente al bit de acarreo que sale del MSB. En otras palabras: ocurre al sumar dos números del mismo signo y obtener un resultado del signo opuesto.
            \end{itemize}
        }



    % - Problema 20
    \vspace{\Pspace}
\item La página 22 de los \href{https://unisonmx.sharepoint.com/sites/Arquitecturadecomputadoras2026-1/Class%20Materials/Slides/09-x86-Overview.pdf#page=22}{apuntes sobre x86} dice que la instrucción \mintinline{mips}{jl} (jump if less) brinca si la bandera de signo es diferente a la de overflow. Explicar la lógica de esto y por qué funciona. \par
        % Respuesta:
        \vspace{\Aspace} \par
        { \color{azul}
            La instrucción \asm{jl} se basa en el resultado de una resta interna ($A - B$). La lógica de saltar si \textbf{SF $\neq$ OF} funciona porque el Overflow Flag actúa como un corrector del Sign Flag.
            \begin{itemize}
                \item \textbf{Caso sin Overflow (OF = 0):} El signo del resultado es verídico. Si $A < B$, la resta arroja un número negativo ($SF = 1$). Como $1 \neq 0$, la condición se cumple y el salto ocurre.
                \item \textbf{Caso con Overflow (OF = 1):} El signo del resultado es erróneo (se ha invertido). Si $A < B$ pero ocurre un desbordamiento, el signo parecerá positivo ($SF = 0$). Sin embargo, como $0 \neq 1$, la instrucción identifica que el resultado ``debió" ser negativo y ejecuta el salto.
            \end{itemize}

            En resumen, si el signo es negativo y todo es normal, o si el signo es positivo pero sabemos que hubo un error de desbordamiento, significa matemáticamente que $A$ es menor que $B$. 
        }
    \end{enumerate}
\end{document}
